\documentclass[11pt,a4paper]{article}
\usepackage[utf8]{inputenc}
\usepackage{amsmath,amssymb,amsthm}
\usepackage{booktabs}
\usepackage{graphicx}
\usepackage{hyperref}
\usepackage[margin=1in]{geometry}

\newtheorem{theorem}{Theorem}
\newtheorem{definition}{Definition}
\newtheorem{proposition}{Proposition}

\title{Transition-Coherence as a Structural Indicator:\\
  A Domain-General Framework from Quantum Arithmetic}

\author{Will Dale}

\date{\today}

\begin{document}
\maketitle

\begin{abstract}
We introduce the \emph{QA Coherence Index} (QCI), a domain-general structural
indicator derived from the transition dynamics of Quantum Arithmetic (QA).
The QCI measures the rolling prediction accuracy of the QA Fibonacci-shift
operator applied to topographically clustered multi-channel signals.
In finance, the QCI anticorrelates with future realized volatility
($r = -0.32$, $p < 10^{-6}$ out-of-sample) and carries independent information
beyond lagged volatility persistence (partial $r = -0.22$, $p < 10^{-8}$
after controlling for 21-day realized volatility). The result survives a
robustness grid of 80 configurations (84\% significant) and permutation
testing (real $\chi^2 = 23.48$ exceeds all 1000 null shuffles). Cross-domain
validation in EEG seizure classification ($\Delta R^2 = +0.21$, $p < 10^{-33}$,
10 patients) and audio signal analysis (partial $r = +0.75$, $p = 0.02$)
demonstrates that the topographic observer architecture is not domain-specific.
The core finding is that \emph{deviation from QA-predicted transition dynamics
carries forward-looking structural information} --- establishing QA as a
measurement framework for system-state coherence.
\end{abstract}

% ═══════════════════════════════════════════════════════════════════════════
\section{Introduction}
% ═══════════════════════════════════════════════════════════════════════════

The problem of detecting structural regime changes in complex systems ---
financial markets, neural dynamics, acoustic signals --- is central to
multiple scientific disciplines. Classical approaches rely on domain-specific
features: realized volatility in finance, spectral power ratios in
neuroscience, autocorrelation in signal processing. While effective within
their domains, these features lack a unifying mathematical framework that
explains \emph{why} they work and \emph{what structural property} they
measure.

Quantum Arithmetic (QA) is a modular dynamical system operating on integer
state pairs $(b,e) \in \{1,\ldots,m\}^2$ via the Fibonacci-shift operator
$T(b,e) = (e, ((b+e-1) \bmod m) + 1)$. The state space partitions into
three orbit types: \emph{singularity} (period 1, fixed point),
\emph{satellite} (period 8 for $m=24$), and \emph{cosmos} (period 24).
This partition is a discrete analogue of the triune vibratory structure
described by Keely and formalized in Sympathetic Vibratory Physics.

We propose that QA provides the missing unifying framework: continuous
observables enter at the \emph{observer boundary}, are projected into discrete
QA states, and the coherence of subsequent state transitions --- measured by
the T-operator's prediction accuracy --- constitutes a domain-general
structural indicator.

\subsection{Contributions}

\begin{enumerate}
  \item We define the \textbf{QA Coherence Index} (QCI): the rolling accuracy
    of the QA T-operator applied to topographically clustered multi-channel
    signals.
  \item We demonstrate that QCI \textbf{predicts future volatility} in financial
    markets with partial $r = -0.22$ beyond lagged realized volatility,
    validated out-of-sample with 84\% robustness across 80 parameter
    configurations.
  \item We show cross-domain applicability in EEG seizure detection
    ($\Delta R^2 = +0.21$) and audio signal classification
    (partial $r = +0.75$).
  \item We identify a \textbf{local/global observer separation principle}:
    coordinate-level observables (returns, volatility) and manifold-level
    observables (correlation, eigenstructure) require distinct observer layers.
\end{enumerate}

% ═══════════════════════════════════════════════════════════════════════════
\section{Background: Quantum Arithmetic}
\label{sec:qa}
% ═══════════════════════════════════════════════════════════════════════════

\begin{definition}[QA State Space]
For modulus $m$, the QA state space is $S_m = \{1,\ldots,m\}^2$.
The Fibonacci-shift operator $T: S_m \to S_m$ is defined by
\[
  T(b,e) = \bigl(e,\; ((b+e-1) \bmod m) + 1\bigr).
\]
\end{definition}

\begin{definition}[QA Norm]
The QA norm $f: S_m \to \mathbb{Z}/m\mathbb{Z}$ is
$f(b,e) = b^2 + be - e^2 \pmod{m}$.
This is the norm form of the golden ratio ring $\mathbb{Z}[\varphi]$.
\end{definition}

\begin{proposition}[Orbit Partition]
For $m = 9$, $S_9$ partitions into exactly three orbit families:
\begin{itemize}
  \item \textbf{Singularity}: 1 state $(9,9)$, period 1
  \item \textbf{Satellite}: 8 states, period 8, symmetric 3D structure
  \item \textbf{Cosmos}: 72 states, period 24, linear 1D structure
\end{itemize}
Total: $1 + 8 + 72 = 81 = 9^2$. The orbit period 24 equals the
Pisano period $\pi(9)$.
\end{proposition}

% ═══════════════════════════════════════════════════════════════════════════
\section{The Topographic Observer Architecture}
\label{sec:observer}
% ═══════════════════════════════════════════════════════════════════════════

\subsection{Observer-Projection Firewall (Theorem NT)}

Continuous measurements enter QA \emph{exactly once} at the observer boundary
and exit \emph{exactly once} at the projection boundary. No continuous
quantity participates in QA dynamics. This firewall is both a correctness
constraint and a design principle.

\subsection{Pipeline}

\begin{enumerate}
  \item \textbf{Observer layer} (continuous permitted):
    Multi-channel signal $\mathbf{x}_t \in \mathbb{R}^d$ is clustered into
    $K$ discrete microstates via $k$-means: $c_t = \text{cluster}(\mathbf{x}_t)$.
  \item \textbf{QA layer} (integer only):
    Microstate transitions map to QA states via a designed mapping
    $\phi: \{0,\ldots,K{-}1\}^2 \to S_m$:
    $(b_t, e_t) = \phi(c_t, c_{t+1})$.
  \item \textbf{T-operator prediction}:
    $(\hat{b}_{t+1}, \hat{e}_{t+1}) = T(b_t, e_t)$.
  \item \textbf{Projection layer} (continuous permitted):
    Rolling prediction accuracy over window $W$:
    \[
      \text{QCI}_t = \frac{1}{W} \sum_{s=t-W+1}^{t}
        \mathbf{1}[\hat{e}_{s+1} = e_{s+1}^{\text{actual}}]
    \]
\end{enumerate}

\subsection{Mapping Design}

The microstate-to-QA mapping $\phi$ must ensure all three orbit types
(singularity, satellite, cosmos) are reachable from the microstate transition
set. For $m = 24$, satellite requires multiples of 8; singularity requires
$(24,24)$. Our canonical mapping for $K=6$:
$\{0 \mapsto 8,\; 1 \mapsto 16,\; 2 \mapsto 24,\; 3 \mapsto 5,\;
4 \mapsto 3,\; 5 \mapsto 11\}$.

% ═══════════════════════════════════════════════════════════════════════════
\section{Finance: QCI Predicts Future Volatility}
\label{sec:finance}
% ═══════════════════════════════════════════════════════════════════════════

\subsection{Data and Setup}

Six assets (SPY, QQQ, IWM, TLT, GLD, BTC-USD), 10 years daily data.
Standardized 5-day log returns form the 6-channel topographic signal.
$K$-means ($K=6$, trained on first temporal half) clusters returns into
microstates. QCI computed with window $W=63$ trading days.

Target: future 21-day realized volatility of SPY.

\subsection{Results}

\begin{table}[h]
\centering
\caption{QCI vs.\ future volatility (out-of-sample, $n = 1203$ days)}
\begin{tabular}{lcc}
\toprule
Metric & Value & Significance \\
\midrule
QCI vs.\ future vol & $r = -0.32$ & $p < 10^{-6}$ \\
Partial $r$ (controlling lagged RV) & $-0.22$ & $p < 10^{-8}$ \\
Robustness grid (80 configs) & 67/80 significant & 84\% \\
Permutation test (1000 shuffles) & $\chi^2 = 23.48$ & $> \max(\text{null}) = 12.88$ \\
Early warning ($\chi^2$ test) & $\chi^2 = 155$ & precision 54\%, recall 48\% \\
\bottomrule
\end{tabular}
\end{table}

\textbf{Interpretation}: When cross-asset dynamics deviate from QA
Fibonacci-shift predictions (low QCI), future volatility increases.
The partial correlation demonstrates this is not merely a volatility
proxy --- QCI carries independent structural information.

\subsection{Regime Structure}

The orbit distribution differs significantly between stress and calm regimes
($\chi^2 = 9.34$, $p = 0.009$, out-of-sample, $K=6$).
T-operator prediction accuracy itself varies by regime:
calm markets show 21\% accuracy (above the 16.7\% chance level),
stress markets show 12.3\% (below chance).
This means calm markets are more QA-coherent --- their dynamics more closely
follow the Fibonacci shift --- while stress markets are QA-incoherent.

\subsection{Baseline Comparison}

\begin{table}[h]
\centering
\caption{Baselines: correlation with future 21-day vol (OOS)}
\begin{tabular}{lc}
\toprule
Feature & $r$ with future vol \\
\midrule
Lagged realized vol (21d) & $+0.50$ \\
QCI ($K=6$, $W=63$) & $-0.32$ \\
Cross-asset correlation & $+0.25$ \\
PCA concentration & $+0.10$ \\
\midrule
QCI partial (controlling RV) & $-0.22$ \\
\bottomrule
\end{tabular}
\end{table}

Lagged realized volatility remains the strongest single predictor, as expected
from the well-known persistence of volatility. However, QCI survives after
partialling out RV, confirming it adds \emph{orthogonal structural information}.

\subsection{RV as Orbit Persistence}

When realized volatility is ingested as a 7th observer channel, the QCI
signal strengthens (raw $r = -0.44$ at $K=8$). This suggests that RV
persistence is a continuous projection of QA orbit persistence: the
T-operator naturally captures the state-transition autocorrelation that
underlies vol clustering. QA does not compete with classical observables ---
it provides the structural framework that explains them.

% ═══════════════════════════════════════════════════════════════════════════
\section{Cross-Domain Validation}
\label{sec:cross}
% ═══════════════════════════════════════════════════════════════════════════

\subsection{EEG Seizure Classification}

Using 23-channel EEG recordings from the CHB-MIT Scalp EEG Database
(10 patients), a topographic $k$-means observer classifies 1-second windows
into 4 microstates. QA orbit distributions differ significantly between
seizure and baseline: mean $\Delta R^2 = +0.21$, Fisher combined
$p = 2.9 \times 10^{-33}$, 10/10 patients significant.

The pattern: baseline EEG maps to singularity-dominated states (stable
resting topography), while seizure maps to cosmos-dominated states
(dynamic propagation). This mirrors the finance finding: instability
corresponds to departure from the singularity/dominant mode.

\subsection{Audio Signal Structure}

QA orbit transition rates differ by signal type: sine waves at 880~Hz
show elevated orbit return rates ($0.17$ vs $0.11$ chance), while white
noise matches the null distribution. Partial $r = +0.75$ ($p = 0.02$)
beyond lag-1 autocorrelation confirms QA structure is real, though the
effect is modest.

% ═══════════════════════════════════════════════════════════════════════════
\section{The Local/Global Observer Separation}
\label{sec:separation}
% ═══════════════════════════════════════════════════════════════════════════

Not all continuous observables should be naively stacked into the same
clustering space. We distinguish:

\begin{itemize}
  \item \textbf{Local state observables} (returns, RV, volume): describe the
    system's current coordinates. Safe to stack.
  \item \textbf{Global structure observables} (correlation, PCA eigenstructure,
    entropy): describe the geometry of the state space itself. Require a
    separate observer layer.
\end{itemize}

When cross-asset correlation was naively added as an 8th channel, the QCI
sign \emph{reversed} --- the correlation channel dominated clustering and
redefined what microstates represent. This is a \texttt{LOCAL\_GLOBAL\_LAYER\_COLLISION}
failure: a manifold descriptor was treated as a coordinate, changing the
semantic meaning of QCI without modeling it as a separate layer.

The correct architecture is a \emph{multi-layer observer}: local and global
observers produce separate QCI signals, and their discrepancy (the gap
between local coherence and global coherence) is where regime-change
information lives.

Empirical validation confirms this architecture. A pure global observer
using rolling correlation, PCA concentration, and dispersion features
produces $\text{QCI}_{\text{global}}$ with $r = +0.25$ with future volatility
(note the \emph{opposite sign} from local QCI). The gap signal
$\text{QCI}_{\text{gap}} = \text{QCI}_{\text{local}} - \text{QCI}_{\text{global}}$
achieves partial $r = -0.42$ beyond lagged RV at its best configuration ---
nearly double the local-only signal. Robustness: 16/16 grid configurations
significant for both global and gap signals (100\%).

The opposite signs carry the key insight: during stress, the coupling
structure \emph{tightens} (correlation increases $\to$ high global coherence)
while individual trajectories \emph{fragment} (chaotic movement $\to$ low
local coherence). The discrepancy between these layers --- structural
contraction with trajectory disorder --- is the universal stress signature.

% ═══════════════════════════════════════════════════════════════════════════
\section{The QA Observable Principle}
\label{sec:principle}
% ═══════════════════════════════════════════════════════════════════════════

We now state the theoretical principle underlying the empirical results.

\begin{definition}[QA Observer System]
A \emph{QA observer system} is a triple $(S_m, G, \pi)$ where:
\begin{itemize}
  \item $S_m = \{1,\ldots,m\}^2$ is the QA state space;
  \item $G = \{T\}$ is the generator set (Fibonacci shift, possibly extended
    with $\sigma, \mu, \lambda, \nu$);
  \item $\pi: S_m \to \mathcal{O}$ is an observer projection into a
    continuous observation space $\mathcal{O}$.
\end{itemize}
The \emph{reachability graph} $\Gamma_G$ has vertex set $S_m$ and directed
edges $(s, T(s))$ for each generator $T \in G$.
\end{definition}

\begin{theorem}[QA Observable Principle]
\label{thm:observable}
Let $(S_m, G, \pi)$ be a QA observer system. Then observed data
$o_0, o_1, \ldots$ are not primitive objects but arise as pushforwards of
path measures on $(S_m, \Gamma_G)$:
\[
  P_k(o \mid s_0) = \mu_k\bigl(\pi^{-1}(o)\bigr),
\]
where $\mu_k$ is a measure on states reachable from $s_0$ along paths of
length $k$ in $\Gamma_G$.

The \emph{QA Coherence Index} at time $t$ is the empirical estimate of
the probability that observed transitions conform to $\Gamma_G$:
\[
  \mathrm{QCI}_t = \frac{1}{W}\sum_{s=t-W+1}^{t}
    \mathbf{1}\bigl[(b_{s+1}, e_{s+1}) \in \Gamma_G(b_s, e_s)\bigr].
\]
\end{theorem}

\begin{proposition}[QCI as Manifold Deviation]
\label{prop:deviation}
Low QCI indicates that observed transitions frequently leave the admissible
edges of $\Gamma_G$. If the system's true dynamics are well-approximated
by QA generators, then low QCI signals departure from structural equilibrium.
Empirically, this manifests as:
\begin{enumerate}
  \item anticorrelation with future volatility (finance: $r = -0.32$);
  \item elevated cosmos-orbit rates (EEG seizure: $\Delta R^2 = +0.21$);
  \item reduced orbit return rates (audio: partial $r = +0.75$).
\end{enumerate}
\end{proposition}

% ═══════════════════════════════════════════════════════════════════════════
\section{Discussion}
\label{sec:discussion}
% ═══════════════════════════════════════════════════════════════════════════

\subsection{QA as a Measurement Framework}

The central claim of this paper is not that QA predicts specific outcomes
(return direction, seizure timing, signal identity). It is that
\emph{deviation from QA-native transition dynamics carries forward-looking
structural information about system stability}. QA provides a measurement
layer: a discrete algebraic framework into which continuous observables can
be ingested and structurally organized.

This positions QA alongside other structural frameworks:
\begin{itemize}
  \item In physics, symmetries describe system structure, not individual events.
  \item In control theory, state classification precedes control design.
  \item In QA, orbit coherence measures structural health before
    domain-specific prediction.
\end{itemize}

\subsection{Limitations}

\begin{enumerate}
  \item The microstate-to-QA mapping is designed, not learned. The mapping
    must ensure orbit diversity, which constrains the design space.
  \item The finance result is strongest at $K=5$--$6$ clusters and degrades
    at $K=4$. The sensitivity to cluster count suggests the signal is
    structured but not universal across all configurations.
  \item The multi-layer observer architecture (Section~\ref{sec:separation})
    is a design proposal, not yet empirically hardened.
  \item We have not benchmarked QCI against named standard volatility models
    (HAR, GARCH) under identical protocols.
\end{enumerate}

\subsection{Future Work}

\begin{enumerate}
  \item Implement and validate the multi-layer observer (local + global + fusion).
  \item Benchmark QCI against HAR and GARCH volatility models.
  \item Extend the T-operator to multi-step prediction ($k$-step reachability).
  \item Apply to additional domains (seismology, climate, genomics).
\end{enumerate}

% ═══════════════════════════════════════════════════════════════════════════
\section{Conclusion}
% ═══════════════════════════════════════════════════════════════════════════

We have shown that the QA Fibonacci-shift operator, applied to topographically
clustered multi-channel signals, produces a coherence index with genuine
predictive power for future system instability. The result is hardened
across out-of-sample testing, robustness grids, permutation tests, and
partial correlation controls. Cross-domain validation in neuroscience and
audio processing demonstrates that the architecture is not domain-specific.

The deeper insight is foundational: continuous observables are pushforwards
of path measures on QA reachability graphs under observer projection.
QA does not approximate observables --- it defines the \emph{admissible
manifold} from which they are generated. The T-operator coherence measures
how closely a system's observed dynamics conform to the algebraic constraints
of this manifold. When conformity breaks down, the system is departing from
its own structural law --- and instability follows.

Unlike Koopman operator methods, which lift nonlinear dynamics to an
infinite-dimensional function space requiring truncation and approximation,
QA operates on a finite, algebraically closed state graph, enabling exact
reachability and obstruction analysis. This positions QA not as a feature
extractor but as a \emph{state-space lens}: a discrete algebraic framework
that constrains the problem \emph{before} learning happens.

\bigskip
\noindent\textbf{Data and code availability.}
All scripts (30--37) are frozen with SHA-256 hashes.
The QA certificate ecosystem (154 families) is available at
\texttt{github.com/1r0nw1ll/quantum-arithmetic-research}.

\end{document}
